\chapter{Complete Implementation of ML-KEM (FIPS 203)}
\chaptermark{ML-KEM (FIPS 203)}
\label{ch:fips203}

\section{Learning objectives}

This chapter assembles the pieces built earlier in the book -- the negacyclic ring and the NTT (Chapter~\ref{ch:ntt}), the toy scheme (Chapter~\ref{ch:toymlkem}), and centered binomial sampling (Chapter~\ref{ch:cbd}) -- into the real standard. FIPS 203, \textbf{ML-KEM} (from the submission CRYSTALS-Kyber)~\cite{fips203,kyber2018}, is the module-lattice-based key-encapsulation standard, and the first of NIST's post-quantum standards to be finalized (2024). All algorithm and data-flow details in this chapter follow FIPS 203~\cite{fips203}.

After finishing this chapter, you should be able to:

\begin{enumerate}
  \item follow the full ML-KEM key generation workflow and its data types at every stage;
  \item explain why key generation is almost entirely deterministic once the initial seed is fixed;
  \item follow encapsulation and explain why a KEM encapsulates a \emph{random} secret rather than encrypting chosen data;
  \item follow decapsulation and explain the Fujisaki--Okamoto ``re-encrypt and compare'' transform with implicit rejection;
  \item distinguish the CPA-secure inner scheme (K-PKE) from the CCA-secure KEM (ML-KEM) built on top of it.
\end{enumerate}

\section{Review: Toy Kyber versus real ML-KEM}

In the toy implementation of Chapter~\ref{ch:toymlkem}, we used the following informal key-generation pattern:

\[
A \xleftarrow{\$} R_q^{k\times k}, \qquad s,e \xleftarrow{\$} \chi, \qquad t = A\cdot s + e.
\]

The output was then a key pair \((pk, sk)\).

The most important difference is that all random objects were generated directly and independently.

In real ML-KEM, this is not the case. The algorithm starts from a short random seed and then deterministically expands the rest of the randomness from it.

\section{The real input in FIPS 203}

The standard key generation procedure does not begin by generating a matrix at random.
It begins by generating a short random seed:

\[
d \in \{0,1\}^{256}.
\]

In other words, the input to the standard algorithm is a 256-bit seed. We can think of the whole key generation algorithm as a deterministic function

\[
\mathrm{KeyGen}(d)
\]

rather than as a random process that samples all objects independently.

This is the first major conceptual shift from toy Kyber to the real standard.

\section{Phase 1: expand the seed}

The first step in the standard is to compute

\[
(\rho, \sigma) = G(d),
\]

where:

\begin{itemize}
  \item $\rho$ is used to generate the matrix $A$;
  \item $\sigma$ is used to generate the secret and error polynomials.
\end{itemize}

In the standard, $G$ is built from SHA3-512. In a teaching implementation, we can first abstract it as a black box:

\begin{figure}[H]
\centering
\begin{tikzpicture}[node distance=9mm and 16mm]
  \node[flownode] (d) {$d$};
  \node[flownode, right=of d] (g) {$G(d)$};
  \node[flownode, above right=6mm and 16mm of g] (rho) {$\rho$};
  \node[flownode, below right=6mm and 16mm of g] (sigma) {$\sigma$};
  \draw[flowarrow] (d) -- (g);
  \draw[flowarrow] (g) -- (rho);
  \draw[flowarrow] (g) -- (sigma);
\end{tikzpicture}
\caption{$G$ expands the single seed $d$ into two independent-looking seeds: $\rho$ drives the matrix, $\sigma$ drives the noise.}
\label{fig:seed-expand-g}
\end{figure}

The pedagogical point is that the entire randomness of the later steps is derived from the initial seed $d$.

\section{Phase 2: generate the matrix}

The matrix is not stored as a random object in the public key. Instead, it is reconstructed from a seed:

\[
A = \mathrm{ExpandA}(\rho).
\]

This is a deterministic map:

\[
\mathrm{ExpandA}: \{0,1\}^{256} \to R_q^{k\times k}.
\]

The same seed always yields the same matrix. This is a crucial compression trick in the standard.

\paragraph{SageMath experiment}

For a teaching implementation, we can first simulate the expansion by using a simple deterministic pseudo-random generator.

\begin{lstlisting}[language=Python]
from random import Random
from sage.all import GF, Matrix, PolynomialRing

q = 17
n = 8
k = 2
R.<x> = PolynomialRing(GF(q))
Rq = R.quotient(x^n + 1)


def expand_A(seed):
    rng = Random(str(seed))
    return Matrix(
        Rq,
        k,
        k,
        [Rq([rng.randrange(q) for _ in range(n)])
         for _ in range(k * k)]
    )
\end{lstlisting}

The important thing is not the exact random source at this stage. The important thing is the interface:

\begin{itemize}
  \item same seed $\Rightarrow$ same matrix;
  \item different seed $\Rightarrow$ different matrix.
\end{itemize}

\section{Phase 3: generate the secret vector}

The secret vector is derived from the CBD distribution. In other words,

\[
s \leftarrow \mathrm{CBD}(\sigma).
\]

This is not uniform sampling. The coefficients are small integers sampled from a centered binomial distribution. Each coefficient of the secret polynomial is therefore a small signed value rather than an arbitrary element of the ring.

The secret vector is therefore a vector of small polynomials:

\[
s = (s_0, \dots, s_{k-1}).
\]

\section{Phase 4: generate the error vector}

The error vector is generated in a very similar way, again using CBD-based sampling:

\[
e \leftarrow \mathrm{CBD}(\sigma).
\]

The exact detail depends on the parameter set, but the conceptual point is the same: the noise is derived from a short seed through a deterministic sampling procedure.

\section{Phase 5: compute the public key in the NTT domain}

FIPS~203 samples the secret and error in coefficient form, then transforms them. \textsc{ExpandA} already produces the matrix in the NTT representation. The public vector is computed as
\[
\hat t=\hat A\circ\hat s+\hat e,
\]
where $\circ$ uses the specified \textsc{MultiplyNTTs} operation: 128 quadratic base products, not unrestricted coordinate-wise scalar multiplication. The PKE public key encodes $(\hat t,\rho)$ and the PKE secret key encodes $\hat s$; an inverse NTT is used later where coefficient-domain values are required. This is the data flow of FIPS~203, Sections~4.3.1 and~5.1.

\section{Phase 7: encode the public key}

The real public key is not the pair \((A,t)\). Instead, it is essentially the pair

\[
(\rho, t).
\]

The reason is that the matrix $A$ can be regenerated from the seed $\rho$. Therefore the public key stores only what is truly necessary:

\begin{itemize}
  \item the seed $\rho$;
  \item the vector $t$.
\end{itemize}

\section{KeyGen data flow}

The whole process can be summarized as follows:

\begin{figure}[H]
\centering
\begin{tikzpicture}[every node/.style={font=\scriptsize}]
  \node[flownode] (d) at (5,9.6) {Random seed $d$};
  \node[flownode] (g) at (5,8.2) {$G(d) = \mathrm{SHA3\text{-}512}(d)$};
  \node[flownode] (rho) at (2.3,6.8) {$\rho$};
  \node[flownode] (sigma) at (7.7,6.8) {$\sigma$};
  \node[flownode] (expandA) at (2.3,5.4) {ExpandA$(\rho)=\hat A$};
  \node[flownode] (cbd) at (7.7,5.4) {CBD$(\sigma)$};
  \node[flownode] (s) at (6.4,4.0) {$s\to\hat s$};
  \node[flownode] (e) at (9.0,4.0) {$e\to\hat e$};
  \node[flownode] (mult) at (5,2.6) {$\hat t = \hat A\circ\hat s + \hat e$};
  \node[keynode] (pk) at (5,1.1) {Public key $(\rho,\hat t)$};

  \draw[flowarrow] (d) -- (g);
  \draw[flowarrow] (g) -- (rho);
  \draw[flowarrow] (g) -- (sigma);
  \draw[flowarrow] (rho) -- (expandA);
  \draw[flowarrow] (sigma) -- (cbd);
  \draw[flowarrow] (cbd) -- (s);
  \draw[flowarrow] (cbd) -- (e);
  \draw[flowarrow] (expandA) -- (mult);
  \draw[flowarrow] (s) -- (mult);
  \draw[flowarrow] (e) -- (mult);
  \draw[flowarrow] (mult) -- (pk);
  \draw[flowarrow] (rho.west) -- ++(-1.3,0) |- (pk.west);
\end{tikzpicture}
\caption{The K-PKE KeyGen data flow. The seed $d$ expands into $\rho$ and the sampling seed; \textsc{ExpandA} produces $\hat A$, while sampled coefficient-domain vectors are transformed before the NTT-domain multiplication. The public key stores $\rho$ and $\hat t$, never $A$ itself.}
\label{fig:keygen-flow}
\end{figure}

This flow should be remembered because it reappears inside encapsulation and decapsulation below. Written as pseudocode, the seven phases above collapse into a single short algorithm:

\begin{algorithm}[H]
\caption{K-PKE.KeyGen (FIPS 203~\cite{fips203}, teaching form)}
\label{alg:kpke-keygen}
\begin{algorithmic}[1]
\Require Random seed $d \in \{0,1\}^{256}$
\Ensure Encoded public key $(\rho,\hat t)$; encoded secret key $\hat s$
\State $(\rho,\sigma) \gets G(d)$ \Comment{Phase 1: seed expansion, $G=\mathrm{SHA3\text{-}512}$}
\State $\hat A \gets \textsc{ExpandA}(\rho)$ \Comment{NTT representation of the deterministic matrix}
\State $s \gets \textsc{CBD}_{\eta_1}(\sigma)$ \Comment{Phase 3: secret vector}
\State $e \gets \textsc{CBD}_{\eta_1}(\sigma)$ \Comment{Phase 4: error vector}
\State $\hat s,\hat e \gets \mathrm{NTT}(s),\mathrm{NTT}(e)$
\State $\hat t \gets \hat A\circ\hat s+\hat e$ \Comment{use \textsc{MultiplyNTTs} for every ring product}
\State \Return $\big((\rho,\hat t),\ \hat s\big)$ \Comment{encode the NTT-domain values as specified}
\end{algorithmic}
\end{algorithm}

Every line here is deterministic given $d$; the only place fresh randomness enters the entire computation is the sampling of $d$ itself.

\section{From K-PKE to ML-KEM: the two layers}

The algorithm above is not quite ML-KEM. It is the key generation of \textbf{K-PKE}, the underlying \emph{public-key encryption} scheme, which is only \emph{CPA-secure} -- safe against a passive eavesdropper, but not against an attacker who submits chosen ciphertexts and watches how decryption responds. ML-KEM is K-PKE wrapped in one more layer. K-PKE provides three routines,
\[
\textsc{K-PKE.KeyGen},\qquad \textsc{K-PKE.Encrypt}(ek, m, r),\qquad \textsc{K-PKE.Decrypt}(dk, c),
\]
and the \textbf{Fujisaki--Okamoto (FO) transform} turns them into a \emph{CCA-secure} key-encapsulation mechanism. ML-KEM key generation simply calls K-PKE key generation and stores a little extra in the secret key:

\begin{algorithm}[H]
\caption{ML-KEM.KeyGen (teaching form)}
\label{alg:mlkem-keygen}
\begin{algorithmic}[1]
\Require Independent random values $d,z \in \{0,1\}^{256}$
\Ensure Encapsulation key $ek$; decapsulation key $dk$
\State $\big((\rho,t),\ dk_{\mathrm{PKE}}\big) \gets \textsc{K-PKE.KeyGen}(d)$ \Comment{Algorithm~\ref{alg:kpke-keygen}}
\State $ek \gets (\rho, t)$
\State \Return $ek,\ \ dk \gets \big(dk_{\mathrm{PKE}},\ ek,\ H(ek),\ z\big)$
\end{algorithmic}
\end{algorithm}

The encapsulation key $ek=(\rho,t)$ is what a KEM publishes. The decapsulation key carries the PKE secret $dk_{\mathrm{PKE}}=s$, a copy of $ek$ and its hash $H(ek)$ (both needed to re-encrypt during decapsulation), and the random value $z$ whose role becomes clear below.

\section{Encapsulation}

Encapsulation produces a fresh \emph{shared secret} $K$ and a ciphertext $c$ from which the holder of $dk$ can recover the same $K$. The design choice that defines a KEM -- and the answer to ``why encapsulate a random secret instead of encrypting a chosen message?'' -- is that the encrypted message is \emph{not chosen by the sender at all}. Encapsulation samples a random $32$-byte value $m$, and derives both the shared key and the encryption randomness by \emph{hashing} it. That the randomness is a deterministic function of the message is exactly what the FO transform will exploit.

\begin{algorithm}[H]
\caption{ML-KEM.Encaps (teaching form)}
\label{alg:mlkem-encaps}
\begin{algorithmic}[1]
\Require Encapsulation key $ek$
\Ensure Shared secret $K$; ciphertext $c$
\State $m \gets \{0,1\}^{256}$ \Comment{fresh random message, not chosen data}
\State $(K, r) \gets G\big(m \,\|\, H(ek)\big)$ \Comment{$G=\mathrm{SHA3\text{-}512}$: shared key $K$ and encryption seed $r$}
\State $c \gets \textsc{K-PKE.Encrypt}(ek, m, r)$
\State \Return $(K, c)$
\end{algorithmic}
\end{algorithm}

Inside $\textsc{K-PKE.Encrypt}$ the message is hidden in the same noisy-linear way as the toy scheme of Chapter~\ref{ch:toymlkem}, now driven by the seed $r$: the matrix is regenerated as $A=\textsc{ExpandA}(\rho)$, small vectors $y,e_1,e_2$ are sampled from CBD using $r$, and the ciphertext is the compressed pair
\[
u = A^{T} y + e_1, \qquad v = t^{T} y + e_2 + \mathrm{Encode}(m).
\]
The shared secret $K$ itself never travels on the wire; only $c$ does, and the recipient reconstructs $K$ from it.

\section{From CPA to CCA: the Fujisaki--Okamoto transform}
\label{sec:fo}

Everything built so far -- the toy scheme of Chapter~\ref{ch:toymlkem}, the real
noise of Chapter~\ref{ch:cbd}, and K-PKE above -- is a \emph{public-key
encryption} scheme secure only against a passive eavesdropper. One more layer is
needed before any of it is safe to deploy, and FIPS 203 is blunt that the layer
is not optional: K-PKE ``is not IND-CCA2-secure and shall not be used as a
stand-alone scheme''~\cite{fips203}. This section explains what that missing
property is and how the \textbf{Fujisaki--Okamoto (FO) transform}~\cite{fo1999}
supplies it. The same transform reappears in every other KEM in this book --
Classic McEliece (Chapter~\ref{ch:mceliece}), HQC (Chapter~\ref{ch:hqc}), and
SIKE (Chapter~\ref{ch:isogeny}) -- so it is worth meeting properly once.

\paragraph{Two attack models.}
Security notions for encryption differ in what the attacker may \emph{do}, not
in what they are trying to learn. Under both notions the attacker picks two
messages, is given an encryption of one of them, and must say which:

\begin{itemize}
  \item \textbf{IND-CPA} -- indistinguishability under \emph{chosen-plaintext}
    attack. The attacker holds the public key and so can encrypt anything they
    wish. This models a passive eavesdropper.
  \item \textbf{IND-CCA2} -- indistinguishability under \emph{adaptive
    chosen-ciphertext} attack. The attacker may additionally submit ciphertexts
    of their own construction to the honest key holder and observe the result,
    both before and after seeing the challenge; only decrypting the challenge
    itself is barred.
\end{itemize}

The distance between the two is the distance between an attacker who
\emph{listens} and one who \emph{interacts}. Real protocols grant the second
kind of access almost by definition: a server decrypts whatever arrives on the
wire, and what it does next -- an error, a timing difference, a retry -- is
visible to whoever sent it.

\paragraph{Why the passive scheme breaks.}
The failure is concrete rather than theoretical. A K-PKE ciphertext is a pair
$(\mathbf{u}, v)$ in which the message bits are carried by $v$, each scaled by
$\lceil q/2 \rceil$. Arithmetic is coordinatewise, so an attacker who intercepts
$(\mathbf{u}, v)$ and forwards $(\mathbf{u},\ v + \lceil q/2 \rceil)$ hands the
receiver a perfectly valid ciphertext that decrypts to the \emph{complement} of
the original message -- without knowing the secret key and without ever learning
the message. This is \textbf{malleability}: ciphertexts can be edited into
related ciphertexts. The linearity that makes lattice encryption fast is exactly
what makes it malleable.

The decryption oracle leaks, too. Feed the key holder a stream of slightly
malformed ciphertexts and record which decrypt cleanly; each answer is one
inequality constraining the secret, and enough of them recover it. This is why a
passively secure lattice scheme must never be exposed directly.

\paragraph{The key idea: make encryption deterministic.}
Malleability survives only because a ciphertext carries no evidence of how it
was built. Encryption is randomized -- $\textsc{Encrypt}(ek, m; r)$ for a fresh
random $r$ -- so a message has \emph{many} valid ciphertexts and the receiver
cannot tell an honest one from a doctored one.

FO removes the freedom. Rather than sampling $r$, \emph{derive} it from the
message: $r = G(m)$. Encryption becomes a function of the message alone, so each
message now has exactly \emph{one} valid ciphertext -- and that makes ciphertexts
checkable. The receiver decrypts to a candidate $m'$, recomputes the ciphertext
that $m'$ \emph{must} have produced, and compares. An honest ciphertext matches;
anything the attacker altered does not.

The KEM setting is what makes this painless. Because a KEM transports a
\emph{random} secret rather than a chosen message (Chapter~\ref{ch:toymlkem}),
deriving the randomness from the message costs nothing -- the message was never
the point. Encapsulation samples a random $m$ and hashes it into both the shared
key and the encryption randomness, which is precisely what
Algorithm~\ref{alg:mlkem-encaps} does with $(K, r) \gets G(m \,\|\, h)$.

\paragraph{The transform in two steps.}
The modern presentation, due to Hofheinz, H\"ovelmanns, and Kiltz~\cite{hhk2017}
and cited by FIPS 203 itself, splits FO into two independent pieces:

\begin{itemize}
  \item the \textbf{$T$ transform} derandomizes, turning an IND-CPA encryption
    scheme into a \emph{deterministic} one whose ciphertexts can be validated by
    re-encryption;
  \item the \textbf{$U$ transform} wraps that deterministic scheme into a KEM,
    deriving the shared secret by hashing and deciding what to do when the
    re-encryption check fails.
\end{itemize}

Separating the two is what lets the security proof be assembled from parts, and
it is why the literature speaks of a \emph{family} of FO variants rather than one
algorithm. ML-KEM uses one specific member of that family, which the next section
spells out.

\section{Decapsulation and the Fujisaki--Okamoto transform}

Decapsulation must recover $K$ from $c$ using $dk$, and -- this is the entire point of the FO transform -- it must do so in a way that gives an attacker \emph{no} usable signal from submitting malformed ciphertexts. The mechanism is \emph{re-encrypt and compare}: decrypt to a candidate message, deterministically re-run encapsulation from it, and check that the ciphertext that \emph{would} have been produced matches the one received.

\begin{algorithm}[H]
\caption{ML-KEM.Decaps (teaching form)}
\label{alg:mlkem-decaps}
\begin{algorithmic}[1]
\Require Decapsulation key $dk=(dk_{\mathrm{PKE}}, ek, h, z)$; ciphertext $c$
\Ensure Shared secret $K$
\State $m' \gets \textsc{K-PKE.Decrypt}(dk_{\mathrm{PKE}}, c)$
\State $(K', r') \gets G(m' \,\|\, h)$
\State $c' \gets \textsc{K-PKE.Encrypt}(ek, m', r')$ \Comment{re-encrypt from the recovered message}
\If{$c' = c$}
  \State \Return $K'$ \Comment{ciphertext was honestly formed}
\Else
  \State \Return $J(z \,\|\, c)$ \Comment{\emph{implicit rejection}: a pseudorandom key from the secret $z$}
\EndIf
\end{algorithmic}
\end{algorithm}

\begin{figure}[H]
\centering
\begin{tikzpicture}[node distance=6mm and 11mm, every node/.style={font=\scriptsize}]
  \node[flownode] (c) {ciphertext $c$};
  \node[flownode, right=of c] (dec) {$m'=\textsc{Decrypt}(dk,c)$};
  \node[flownode, right=of dec] (g) {$(K',r')=G(m'\,\|\,h)$};
  \node[flownode, below=of g] (re) {$c'=\textsc{Encrypt}(ek,m',r')$};
  \node[diamond, draw=pqcorange, thick, fill=white, align=center, inner sep=1pt, left=of re] (cmp) {$c'=c$?};
  \node[keynode, left=of cmp, align=center] (out) {$K'$ on match\\[-2pt]\scriptsize $J(z\,\|\,c)$ on mismatch};
  \draw[flowarrow] (c) -- (dec);
  \draw[flowarrow] (dec) -- (g);
  \draw[flowarrow] (g) -- (re);
  \draw[flowarrow] (re) -- (cmp);
  \draw[flowarrow] (cmp) -- (out);
\end{tikzpicture}
\caption{ML-KEM decapsulation and the Fujisaki--Okamoto transform. The candidate message is decrypted, encapsulation is re-run deterministically from it, and the recomputed ciphertext is compared with the received one in constant time. A match returns the real shared key; a mismatch returns a pseudorandom key derived from the secret $z$, so success and failure are indistinguishable.}
\label{fig:mlkem-fo}
\end{figure}

The subtlety is the \texttt{else} branch. A naive KEM would return an \emph{error} on a bad ciphertext -- but that error is an oracle: by observing whether decapsulation succeeds, an attacker can mount a chosen-ciphertext attack and, over many queries, extract the secret. \textbf{Implicit rejection} closes this door: on a mismatch, ML-KEM returns a \emph{pseudorandom} key $J(z\,\|\,c)$ derived from the secret $z$, of exactly the same form as a real key. Success and failure become indistinguishable to the caller, so there is no oracle. This is what lifts K-PKE from CPA to \emph{IND-CCA2} security; it is why the comparison $c'=c$ must run in constant time (Chapter~\ref{ch:sidechannel}); and it is the same mechanism the code-based KEM HQC uses (Chapter~\ref{ch:hqc}).

\paragraph{Explicit versus implicit rejection.}
The two ways of handling a failed comparison name the two branches of the FO
family. \emph{Explicit} rejection returns a distinguished error symbol $\bot$;
\emph{implicit} rejection returns a pseudorandom key, as ML-KEM does. Implicit
rejection is the more conservative choice on three counts: it removes the error
signal outright instead of trusting every caller to handle $\bot$ safely; it
keeps decapsulation a single straight-line computation, which is far easier to
make constant-time; and it is the variant for which the tightest security proofs
against quantum adversaries are known (Section~\ref{sec:qrom}).

One detail of ML-KEM's particular variant is easy to miss. The hash $h = H(ek)$
of the encapsulation key is fed into $G$ alongside the message, binding the
shared secret to the public key it was produced under. FIPS 203 also departs
here from the third-round Kyber submission: encapsulation no longer hashes the
ciphertext into the shared secret~\cite{fips203}. That change is safe precisely
because the re-encryption check already ties the ciphertext to the message.

\section{Decryption failures}

Re-encrypt-and-compare assumes decryption returns the message that was
encrypted. For lattice schemes that is only \emph{almost} always true. The noise
of Chapter~\ref{ch:cbd} can, with small probability, push a coefficient past a
rounding boundary and flip a bit of $m'$; the recomputed ciphertext then fails to
match and decapsulation returns the wrong key. This is a correctness failure,
not an attack, and FIPS 203 tabulates its probability~\cite{fips203}:

\begin{table}[H]
\centering
\begin{tabular}{ll}
\toprule
\textbf{Parameter set} & \textbf{Decapsulation failure rate} \\
\midrule
ML-KEM-512  & $2^{-138.8}$ \\
ML-KEM-768  & $2^{-164.8}$ \\
ML-KEM-1024 & $2^{-174.8}$ \\
\bottomrule
\end{tabular}
\caption{Decapsulation failure rates for ML-KEM, from Table 1 of FIPS 203. The
figures hold under the same random-oracle heuristic discussed in
Section~\ref{sec:qrom}.}
\label{tab:mlkem-failure}
\end{table}

These numbers are not merely about reliability. A decryption failure is a
\emph{leak}: it happens only for ciphertexts whose noise correlates with the
secret key, so an attacker able to find failing ciphertexts learns something
about $\mathbf{s}$, and the FO security bound degrades as the failure rate
climbs. That is why the rate is driven far below anything a user would ever
notice, and why parameter selection treats it as a security quantity rather than
an engineering tolerance.

\section{The random oracle model, and its quantum version}
\label{sec:qrom}

The FO proof rests on an assumption worth stating plainly, not least because
FIPS 203 states it too: the analysis holds ``under the heuristic assumption that
hash functions and XOFs behave like uniformly random functions''~\cite{fips203}.
This is the \textbf{random oracle model} (ROM). In it, $G$, $H$, and $J$ are
replaced by genuinely random functions that every party -- honest or hostile --
can only query as a black box. The model is a heuristic, and knowingly so:
SHA3-512 is a fixed public algorithm, not a random function. But proofs in it
have been a reliable guide in practice, and no break of a deployed scheme has
yet been traced to the gap.

For post-quantum cryptography the heuristic needs a second look, and here the
story becomes specific to this book's subject. A quantum attacker does not query
SHA3 as a black box. The function is public, so the attacker can \emph{build the
circuit} and evaluate it on a superposition of inputs, extracting information
about many outputs at once. Modelling that faithfully requires the
\textbf{quantum random oracle model} (QROM), introduced by Boneh and
co-authors~\cite{qrom2011}, in which oracle queries may themselves be
superpositions.

The distinction is not pedantry, because classical ROM proofs do not
automatically survive the move. A great many of them work by \emph{watching} the
attacker's queries and reading off the useful one -- the reduction simply
observes which $m$ the attacker hashed. Against superposition queries that step
is illegitimate: measuring a quantum query disturbs it. Re-establishing FO
security in the QROM was therefore a research programme in its own right,
beginning with the modular analysis of Hofheinz, H\"ovelmanns, and
Kiltz~\cite{hhk2017} and continuing through a line of increasingly tight bounds
for the implicit-rejection variant~\cite{sxy2018}.

This is the honest reason FIPS 203 claims only that ML-KEM ``is believed to
satisfy so-called IND-CCA2 security''~\cite{fips203}. The lattice assumption is
one source of doubt; the random oracle heuristic, in its quantum form, is a
second and quite independent one. If a reader takes one thing from this section,
it should be that ``provably secure'' always carries a model with it -- and that
for every KEM in this book, that model includes a hash function idealized in a
way no real hash function is.

\section{Parameter sets}
\begin{table}[H]
\centering
\begin{tabular}{lccccc}
\toprule
Scheme & $k$ & $\eta_1$ & $(d_u,d_v)$ & NIST level & Comparable to \\
\midrule
ML-KEM-512  & 2 & 3 & $(10,4)$ & 1 & AES-128 \\
ML-KEM-768  & 3 & 2 & $(10,4)$ & 3 & AES-192 \\
ML-KEM-1024 & 4 & 2 & $(11,5)$ & 5 & AES-256 \\
\bottomrule
\end{tabular}
\caption{The three ML-KEM parameter sets share $n=256$ and $q=3329$, but differ in their full parameter tuples. ML-KEM-768 is a common default.}
\label{tab:mlkem-params}
\end{table}

\section{SageMath experiments}

The most useful way to study this chapter is through three short experiments.

\paragraph{Experiment 1: verify ExpandA}

Use the same seed twice:

\begin{lstlisting}[language=Python]
A1 = expand_A(1234)
A2 = expand_A(1234)
\end{lstlisting}

Then verify that \(A_1 = A_2\). Now change the seed and observe that the matrix changes as well.

\paragraph{Experiment 2: inspect CBD statistics}

Generate many samples from the CBD distribution and check that the average is close to zero.

\begin{lstlisting}[language=Python]
from random import randint

def cbd(eta):
    return sum(randint(0, 1) for _ in range(eta)) - \
           sum(randint(0, 1) for _ in range(eta))

samples = [cbd(2) for _ in range(10000)]
print(sum(samples) / len(samples))
\end{lstlisting}

Observe how the variance changes when \(\eta\) changes.

\paragraph{Experiment 3: implement a teaching version of key generation}

Finally, write a function

\begin{lstlisting}
pk, sk = keygen(seed)
\end{lstlisting}

and test the following properties:

\begin{itemize}
  \item same seed gives the same key pair;
  \item different seeds give different key pairs.
\end{itemize}

This gives a first impression of how modern public-key cryptography is built from deterministic pseudo-random procedures.

\paragraph{Experiment 4: one \textsc{MultiplyNTTs} base product}

The NTT-domain multiplication in FIPS~203 is not 256 independent scalar
products. Each coordinate pair is multiplied modulo a quadratic relation. The
following check implements one such pair; it tests the algebraic base case, not
the full standardized transform or its serialization.

\begin{lstlisting}[language=Python]
q = 3329

def multiply_ntt_pair(a, b, gamma):
    a0, a1 = a
    b0, b1 = b
    return ((a0*b0 + gamma*a1*b1) % q,
            (a0*b1 + a1*b0) % q)

# (1 + 2X)(3 + 4X) modulo X^2 - gamma, with gamma = 17.
got = multiply_ntt_pair((1, 2), (3, 4), 17)
assert got == ((3 + 17*8) % q, (4 + 6) % q)
print("one ML-KEM quadratic base product:", got)
\end{lstlisting}

This is precisely the local multiplication used 128 times by
\textsc{MultiplyNTTs}; a conforming ML-KEM implementation must additionally
use FIPS~203's transform ordering, twiddle factors, reductions, encodings, and
test vectors.

\section{A complete SageMath implementation}
\label{sec:fips203-impl}

The experiments above illustrate individual pieces. The companion file
\sagefile{fips203\_mlkem.sage} assembles all of them into a complete,
byte-exact implementation of FIPS~203: every one of the standard's
twenty-one numbered algorithms appears as its own function, named after the
standard and annotated with its algorithm number, for all three parameter
sets.

The algebraic objects are genuine Sage objects rather than opaque byte
arrays. The coefficient domain is the quotient ring itself,
\[
R_q = \mathbb{Z}_q[X]/(X^{256}+1),
\]
and the NTT domain $T_q$ is carried as a Sage vector over $\mathbb{F}_q$:

\begin{lstlisting}[language=Python]
MLKEM_n, MLKEM_q, MLKEM_zeta = 256, 3329, 17

Fq = GF(MLKEM_q)
Rx = PolynomialRing(Fq, 'X')
X  = Rx.gen()
Rq = Rx.quotient(X^MLKEM_n + 1, 'x')   # coefficient domain
Tq = VectorSpace(Fq, MLKEM_n)          # NTT domain

ZETA  = [Fq(MLKEM_zeta)^BitRev7(i) for i in range(128)]
GAMMA = [Fq(MLKEM_zeta)^(2*BitRev7(i) + 1) for i in range(128)]
\end{lstlisting}

With that in place, the transform is the standard's butterfly network written
directly over $\mathbb{F}_q$:

\begin{lstlisting}[language=Python]
def NTT(f):
    """Algorithm 9. Map f in Rq to its NTT representation in Tq."""
    fh = _coeffs(f)
    i, length = 1, 128
    while length >= 2:
        for start in range(0, MLKEM_n, 2*length):
            z = ZETA[i]; i += 1
            for j in range(start, start + length):
                t = z * fh[j + length]
                fh[j + length] = fh[j] - t
                fh[j]          = fh[j] + t
        length //= 2
    return vector(Fq, fh)
\end{lstlisting}

Key generation is then the seven phases of this chapter, line for line:

\begin{lstlisting}[language=Python]
def KPKE_KeyGen(self, d):
    """Algorithm 13. K-PKE.KeyGen(d) -> (ek_PKE, dk_PKE)."""
    k = self.k
    rho, sigma = _G(bytes(d) + bytes([k]))   # domain separator k
    N = 0
    Ah = self._ExpandA(rho)
    s = []
    for _ in range(k):
        s.append(SamplePolyCBD(self.eta1, _PRF(self.eta1, sigma, N))); N += 1
    e = []
    for _ in range(k):
        e.append(SamplePolyCBD(self.eta1, _PRF(self.eta1, sigma, N))); N += 1
    sh = [NTT(si) for si in s]
    eh = [NTT(ei) for ei in e]
    th = [ntt_vec_dot(Ah[i], sh) + eh[i] for i in range(k)]
    ek = b"".join(ByteEncode(12, _ints(t)) for t in th) + rho
    dk = b"".join(ByteEncode(12, _ints(si)) for si in sh)
    return ek, dk
\end{lstlisting}

Note the byte \texttt{k} appended to the seed in the very first line: FIPS~203
domain-separates the three parameter sets so that a seed expanded under the
wrong one yields an unrelated key.

\paragraph{Why a computer-algebra system earns its place here.}
Because $R_q$ is a real Sage ring, the hand-rolled transform can be checked
against the algebra it claims to implement rather than merely against itself.
The function \texttt{verify\_ntt\_is\_a\_ring\_isomorphism()}, defined at the
end of \sagefile{fips203\_mlkem.sage}, confirms three
things using Sage's own arithmetic as the oracle: that $\mathrm{NTT}^{-1}$
inverts $\mathrm{NTT}$; that the $i$-th output coefficient pair really is
$f \bmod (X^2-\gamma_i)$, so the transform is the Chinese-remainder map onto
the $128$ quadratic factors of $X^{256}+1$; and that
\[
\mathrm{NTT}(f\cdot g) \;=\; \textsc{MultiplyNTTs}\big(\mathrm{NTT}(f),\,\mathrm{NTT}(g)\big),
\]
where the left-hand product is computed by Sage in $R_q$. That last identity
is precisely the statement that the NTT is a ring isomorphism $R_q \to T_q$ --
the fact the single-base-product check earlier in this chapter could only
gesture at.

\paragraph{Validation.}
Correctness is not argued, it is measured. The test runner
\sagefile{test\_kat.sage} checks the implementation against NIST's ACVP
test vectors~\cite{acvp}, which carry both the inputs and the expected outputs of every
case: key generation, encapsulation, decapsulation, and the encapsulation-
and decapsulation-key validity checks of Sections~7.2 and~7.3. All outputs
match byte for byte across ML-KEM-512, 768, and 1024. From the book's
\texttt{latex/} directory, \texttt{make kat} runs it.

\section{Pitfalls}

\paragraph{Pitfall 1: thinking that the matrix $A$ is stored explicitly}

This is not true in the standard. The matrix is derived from $\rho$ and therefore regenerated when needed.

\paragraph{Pitfall 2: thinking that CBD is truly random sampling}

CBD is not random sampling in the naive sense. It is a deterministic sampling procedure driven by a seed and an expansion function.

\paragraph{Pitfall 3: thinking that KeyGen calls the random generator at every step}

In the real algorithm, the initial seed $d$ is the only place where fresh randomness is introduced. Everything after that is a deterministic transform.

\paragraph{Pitfall 4: returning an error on decryption failure}

An implementation that signals ``bad ciphertext'' instead of returning the implicit-rejection key $J(z\,\|\,c)$ reintroduces a decryption-failure oracle and breaks CCA security. Both the comparison $c'=c$ and the choice between $K'$ and $J(z\,\|\,c)$ must be constant-time.

\section{Summary}

ML-KEM is a two-layer construction, and this chapter built it end to end:

\begin{itemize}
  \item the inner \textbf{K-PKE} is a CPA-secure public-key encryption whose entire key generation is a deterministic expansion of one seed,
  \[
  d \longrightarrow (\rho,\sigma) \longrightarrow (A,s,e) \longrightarrow t \longrightarrow (\rho,t),
  \]
  publishing only $(\rho,t)$ and never $A$;
  \item \textbf{encapsulation} hides a \emph{random} message $m$, deriving both the shared key $K$ and the encryption randomness $r$ by hashing $m$, so the ciphertext is a deterministic function of $m$;
  \item \textbf{decapsulation} decrypts, re-encrypts, and compares -- returning the real key on a match and a pseudorandom key $J(z\,\|\,c)$ on a mismatch. This is the \textbf{Fujisaki--Okamoto} transform with implicit rejection, and it is what lifts CPA security to \emph{IND-CCA2}.
\end{itemize}

That completes the first NIST post-quantum standard, FIPS 203, from a single random seed to a CCA-secure shared key. The next part turns from key establishment to the other half of public-key cryptography -- digital signatures -- beginning with the lattice-based standards ML-DSA and FN-DSA, which reuse much of the machinery (module lattices, CBD sampling, the NTT) built here.
